\chapter{二次型}\label{ch:b2:quadratic}

\omterm{def:b2:quadratic:def}{二次型}是几何的代数影子：符号为零的部分平坦，
为正的部分向一侧弯曲，为负的
部分向另一侧弯曲。本章将每个实\omterm{def:b2:quadratic:def}{二次型}化为
$\pm$ 平方和（高斯），证明各符号的个数是内在的
（西尔维斯特），再以
\emph{谱定理}为欧氏几何加冕：\omterm{def:b2:quadratic:adjoint}{对称自同态}在
标准正交基中\omterm{def:b2:reduction:diag}{可对角化}——这是应用线性代数
中使用最频繁的一条定理。

\section{双线性型与二次型}

\begin{definition}\label{def:b2:quadratic:def}
实向量空间 $E$ 上的一个\emph{对称双线性型}是
双线性映射 $\varphi \colon E \times E \to \R$，满足 $\varphi(x, y) =
\varphi(y, x)$；与之相伴的
\emph{二次型}\index{二次型}为 $q(x) = \varphi(x, x)$。
$\varphi$ 可由 $q$ 经
\emph{极化恒等式}\index{极化恒等式}：
\[
\varphi(x, y) = \tfrac12\bigl(q(x + y) - q(x) - q(y)\bigr).
\]
在基 $(e_i)$ 下，$\varphi$ 的\emph{矩阵}是对称阵
$B = (\varphi(e_i, e_j))$，且 $q(x) = X^{\mathsf T} B X$；换基矩阵为
$P$ 时，$B$ 被替换为 $P^{\mathsf T} B P$
（\emph{合同}\index{合同}——而非相似！）。
$q$ 的\emph{秩}为 $\operatorname{rk} B$（
不变量：合同即乘以可逆矩阵）。
\end{definition}

\begin{example}
在 $\R^2$ 上：$q(x, y) = x^2 + 4xy + y^2$ 的矩阵为 $\begin{pmatrix}
1 & 2\\ 2 & 1\end{pmatrix}$。内积恰是
其\omterm{def:b2:quadratic:def}{二次型}为正定的对称双线性型；
本章研究一般的、符号不定的情形。
\end{example}

\begin{example}[合同的作用]\label{ex:b2:quadratic:congruence}
取 $q(x, y) = x^2 + 4xy + y^2$（矩阵 $B = \begin{pmatrix} 1
& 2\\ 2 & 1\end{pmatrix}$），新基为 $e_1' = (1, 1)$，
$e_2' = (1, -1)$，即 $P = \begin{pmatrix} 1 & 1\\ 1 &
-1\end{pmatrix}$。于是
\[
P^{\mathsf T}BP
= \begin{pmatrix} 1 & 1\\ 1 & -1\end{pmatrix}
\begin{pmatrix} 1 & 2\\ 2 & 1\end{pmatrix}
\begin{pmatrix} 1 & 1\\ 1 & -1\end{pmatrix}
= \begin{pmatrix} 6 & 0\\ 0 & -2\end{pmatrix} :
\]
在新基下的坐标 $(u, v)$ 中，$q = 6u^2 -
2v^2$——验证：$x = u + v$，$y = u - v$ 代入得 $x^2 + 4xy +
y^2 = 6u^2 - 2v^2$，直接成立。注意新的对角元 $6,
-2$ \emph{并非} $3, -1$ 这两个 $B$ 的\omterm{def:b2:reduction:eigen}{特征值}：\omterm{def:b2:quadratic:def}{合同}
改变尺度，只有相似才保持谱——但符号
一致，正如西尔维斯特定理所要求的。（这里的基
是正交的而非标准正交的；若按
$\frac{1}{\sqrt2}$ 归一化，则对角元除以 $2$，
便恢复出\omterm{def:b2:reduction:eigen}{特征值}。）
\end{example}

\begin{example}[格拉姆行列式度量面积]\label{ex:b2:quadratic:gramarea}
设 $v_1, v_2$ 为欧氏空间中的向量，则格拉姆矩阵 $G =
\bigl(\langle v_i, v_j\rangle\bigr)$ 同时记录长度与
夹角；其\omterm{def:b2:linalg:det}{行列式}则给出\emph{面积}：
\[
\det G = \norm{v_1}^2\norm{v_2}^2 - \langle v_1, v_2\rangle^2
= \norm{v_1}^2\norm{v_2}^2\bigl(1 - \cos^2\theta\bigr)
= \bigl(\norm{v_1}\,\norm{v_2}\sin\theta\bigr)^2 ,
\]
即 $v_1, v_2$ 所张平行四边形面积的平方——
而柯西–施瓦茨不等式恰好就是 $\det G \geq 0$。
例如，取 $v_1 = (1, 2, 2)$，$v_2 = (2, 1, -2)$ 位于
$\R^3$ 中：
\[
G = \begin{pmatrix} 9 & 0\\ 0 & 9 \end{pmatrix},
\qquad
\det G = 81 :
\]
两向量正交且长为 $3$，张成一个
平行四边形（此处为正方形），面积为 $\sqrt{81} = 9$。
结论：既未用到叉积，也未借助 $3$ 维的特殊技巧，
而 $\sqrt{\det G}$ 度量 $k$ 维体积——
在\emph{任意}维数下皆然，是周末问题
第一部分的出发点，也是本卷后文中的
曲面面积积分的基础。
\end{example}

\section{高斯化简与西尔维斯特惯性定律}

\begin{theorem}[高斯化简]\label{thm:b2:quadratic:gauss}
有限维实向量空间上的每个\omterm{def:b2:quadratic:def}{二次型} $q$ 都可以
写成\index{高斯化简}
\[
q = \sum_{i=1}^{s} \ell_i^2 - \sum_{j=1}^{t} m_j^2 ,
\]
其中 $\ell_1, \dots, \ell_s, m_1, \dots, m_t$ 为
线性无关的线性型；等价地，存在某组基使
$q$ 的矩阵为对角阵，对角元为 $+1$（$s$ 次）、$-1$（$t$
次）、$0$。
\end{theorem}

\begin{proof}
对变量个数作归纳，在坐标下考虑 $q(x_1,
\dots, x_n)$。

\emph{情形 1：出现平方项}，不妨设 $a$ 是
$x_1^2$ 的系数，它非零。将所有含 $x_1$ 的项
归并并配成平方：
\[
q = a\Bigl(x_1 + \frac{1}{a}\,\lambda(x_2, \dots,
x_n)\Bigr)^{\!2} + q_1(x_2, \dots, x_n),
\]
其中 $\lambda$ 为线性的，$q_1$ 为其余变量的\omterm{def:b2:quadratic:def}{二次型}：
于是分出一个独立的线性型（它含 $x_1$，其余
变量均不含），对 $q_1$ 应用归纳假设，$\pm$ 号
由 $a$ 的符号经 $\sqrt{\abs a}$ 重新缩放后决定。

\emph{情形 2：无平方项，但有交叉项}，不妨设为 $b\,x_1x_2$，其中
$b \neq 0$。利用恒等式
\[
x_1x_2 = \tfrac14\bigl((x_1 + x_2)^2 - (x_1 - x_2)^2\bigr)
\]
分组之后，令 $q = b\,x_1x_2 + x_1\alpha + x_2\beta +
q_2$（$\alpha, \beta, q_2$ 为其余变量的函数），可直接验证
\[
q = \frac{b}{4}\Bigl[\Bigl(x_1 + x_2 + \frac{\alpha +
\beta}{b}\Bigr)^{2} - \Bigl(x_1 - x_2 + \frac{\beta -
\alpha}{b}\Bigr)^{2}\Bigr] + \widetilde q ,
\]
其中 $\widetilde q$ 不含 $x_1, x_2$：分出两个独立的线性型，
再由归纳法即完成证明。

已收集各型的独立性：将各批按产生的次序排列。
第一批中的型含 $x_1$（情形 1），
或含 $x_1, x_2$（情形 2）；此后各批的型均不含这些
变量。设所收集各型的某一线性组合
为零。取 $x_1$（及 $x_2$）的系数来看，只有
第一批有贡献，而该批中的一个或两个型
显然独立（$\ell$ 本身；或 $\ell \pm
m$，其中 $\ell, m$ 独立）：第一批的系数全为零。
去掉该批并重复此过程：对批次归纳，
所有系数均为零——整族型线性无关，
三角结构得以显式呈现。
\end{proof}

\begin{theorem}[西尔维斯特惯性定律]\label{thm:b2:quadratic:sylvester}
在\cref{thm:b2:quadratic:gauss}中出现的数对 $(s, t)$ 仅依赖于
$q$，与约化过程无关：它是
\emph{符号差}\index{符号差}（$q$ 的）。
此外
\[
s = \max\{\dim F : q|_F \text{ 正定}\},
\]
对 $t$ 同理。
\end{theorem}

\begin{proof}
设 $q = \sum_{i \leq s}\ell_i^2 - \sum_{j\leq t} m_j^2$，并令
$F_+$ 为由（前）对偶向量张成的空间，在这些向量上 $(\ell_i)$
限制为坐标——具体地说：将
独立族 $(\ell_1, \dots, \ell_s, m_1, \dots, m_t)$
扩充为对偶空间 $E^*$ 的一组基，并设 $(u_1, \dots, u_n)$ 为
$E$ 的基，其坐标型恰为这些型（
即前\omterm{def:b2:linalg:dual}{对偶基}：$\ell_i(u_k) = \delta_{ik}$（$k \leq s$），
后面的型在前面的向量上为零）。令 $F_+ =
\operatorname{Vect}(u_1, \dots, u_s)$：对 $x = \sum_{i\leq
s}x_iu_i \in F_+$，有
\[
\ell_i(x) = x_i,
\qquad m_j(x) = 0,
\qquad\text{故}\qquad
q(x) = \sum_{i\leq s}x_i^2 > 0 \quad (x \neq 0) :
\]
$q|_{F_+}$ 正定，且上式中的最大值
$\geq s$。反之，设 $F$ 为任一使 $q|_F$
正定的子空间，记 $G = \{x : \ell_1(x) = \dots = \ell_s(x) =
0\}$，余维数 $\leq s$；在 $G$ 上，$q(x) = -\sum m_j^2 \leq 0$。
于是 $F \cap G = \{0\}$（若其中有非零向量，则 $q > 0$
且 $q \leq 0$），故 $\dim F \leq \dim E - \dim G \leq s$。因此
对\emph{每一种}约化，最大值都等于 $s$：$s$ 是内在的，而 $t
= \operatorname{rk} q - s$ 同样。
\end{proof}

\begin{example}\label{ex:b2:quadratic:gaussexample}
$q(x, y, z) = xy + yz + zx$（无平方项）。取 $x_1 = x$，$x_2 =
y$，则 $q = xy + z(x + y)$，由双平方恒等式得
\[
q = \tfrac14(x + y + 2z)^2 - \tfrac14(x - y)^2 - z^2 ,
\]
（展开即可检验）。三个独立型，其\omterm{thm:b2:quadratic:sylvester}{符号差}为 $(1, 2)$，
秩为 $3$。一个正方向，两个负方向：该型的“光锥”
几何。
\end{example}

\begin{example}[退化型，完全约化]\label{ex:b2:quadratic:degenerate}
$q(x, y, z) = xy + yz$，定义在 $\R^3$ 上：无平方项，故属情形 2，取
分组 $q = y(x + z)$。对独立型的乘积
$y$ 与 $x + z$，应用双平方恒等式：
\[
q = \frac14\bigl(y + x + z\bigr)^2 -
\frac14\bigl(y - x - z\bigr)^2 .
\]
两个线性型 $y + x + z$ 与 $y - x - z$ 是
独立的（它们的差为 $2(x+z)$，和为 $2y$），因此
由西尔维斯特惯性定律直接读出：\omterm{thm:b2:quadratic:sylvester}{符号差}为 $(1, 1)$，秩为 $2$——
\emph{退化}。相伴双线性型的核由
求解 $\varphi(v, \cdot) = 0$ 求得，其矩阵为
$\frac12\begin{pmatrix} 0&1&0\\ 1&0&1\\ 0&1&0\end{pmatrix}$,
核为 $\{y = 0,\ x + z = 0\} = \R\,(1, 0, -1)$，即
$q$ 完全不起作用的方向。小结：秩亏
在\omterm{thm:b2:quadratic:gauss}{高斯化简}中表现为“变量用尽”
——约化只得到三维中的两个平方，
而缺失的那一维恰是核。
\end{example}

\begin{example}[同一二次型，求符号差的两条途径]\label{ex:b2:quadratic:tworoads}
$q(x, y, z) = 2x^2 + 2y^2 + 2z^2 + 2xy + 2yz$，其矩阵为
$\begin{pmatrix} 2 & 1 & 0\\ 1 & 2 & 1\\ 0 & 1 & 2
\end{pmatrix}$。\emph{途径 1，\omterm{thm:b2:quadratic:gauss}{高斯化简}}，依次
配方，
\[
q = 2\Bigl(x + \frac y2\Bigr)^{\!2} + \frac32 y^2 + 2yz + 2z^2
= 2\Bigl(x + \frac y2\Bigr)^{\!2}
+ \frac32\Bigl(y + \frac{2z}{3}\Bigr)^{\!2} + \frac43 z^2 :
\]
三个正平方作用于独立型上，\omterm{thm:b2:quadratic:sylvester}{符号差}为 $(3,
0)$：正定。\emph{途径 2，\omterm{def:b2:reduction:eigen}{特征值}}，矩阵
为三对角矩阵 $2I + N$，其中 $N$ 为邻接矩阵；其
\omterm{def:b2:reduction:eigen}{特征值}为 $2 + \sqrt2$，$2$，$2 - \sqrt2$（可检验
\omterm{def:b2:reduction:eigen}{特征向量} $(1, \pm\sqrt2, 1)$ 与 $(1, 0, -1)$），三者
均为正，结论相同，由
\cref{cor:b2:quadratic:principalaxes}可知。\omterm{thm:b2:quadratic:gauss}{高斯化简}更快；
\omterm{def:b2:reduction:eigen}{特征值}提供的信息更多（它们给出主轴，以及
$q$ 在球面上的极值）。小结：
\omterm{thm:b2:quadratic:gauss}{高斯化简}中的正主元 $2, \frac32, \frac43$ 恰为
顺序主子式的比值 $\frac{\Delta_k}{\Delta_{k-1}}$，
（$\Delta_1 = 2$，$\Delta_2 = 3$，$\Delta_3 = 4$）——
周末问题对一般情形证明了这一点。
\end{example}

\begin{method}[计算符号差：三种途径]\label{met:b2:quadratic:signature}
\begin{enumerate}
  \item \emph{\omterm{thm:b2:quadratic:gauss}{高斯化简}}（总是可行，手算最快）：依次配方，
        无平方项可配时用情形 2；
        统计符号。检验所得的线性型
        线性无关——线性型个数少于变量则意味着有
        核（见\cref{ex:b2:quadratic:degenerate}）。
  \item \emph{顺序主子式}（用于判别正定性）：
        所有 $\Delta_k > 0$ 当且仅当正定
        （见\cref{exo:b2:quadratic:8}）；主元
        $\Delta_k/\Delta_{k-1}$ 甚至给出高斯
        系数（周末问题）。若存在某个
        $\Delta_k = 0$，则此法悄然失效：改用途径 1。
  \item \emph{\omterm{def:b2:reduction:eigen}{特征值}}（信息量最大，计算代价最高）：
        谱的符号
        （见\cref{cor:b2:quadratic:principalaxes}）；还给出主轴，
        以及 $q$ 在
        单位球面上的极值。若无论如何都需要特征结构，
        则优先采用此法。
\end{enumerate}
\end{method}

\section{谱定理}

现设 $E$ 为欧氏空间（内积 $\langle\cdot,\cdot\rangle$，
见本科一年级分册）。

\begin{definition}[伴随；对称自同态]\label{def:b2:quadratic:adjoint}
对 $u \in \mathcal{L}(E)$，其\emph{伴随}\index{伴随}
$u^*$ 是唯一满足下式的自同态：
\[
\langle u(x), y\rangle = \langle x, u^*(y)\rangle
\qquad (x, y \in E);
\]
在标准正交基下，$\operatorname{Mat}(u^*) =
\operatorname{Mat}(u)^{\mathsf T}$。$u$ 称为
\emph{对称自同态}\index{对称自同态}（自伴）
当 $u^* = u$ 时——等价地，其在标准正交基下的矩阵
对称。
\end{definition}

\begin{proof}[伴随的存在与唯一性]
固定 $y$，$x \mapsto \langle u(x), y\rangle$ 是一个
线性型，故（有限维）可写为 $\langle x,
z_y\rangle$ 的形式，其中 $z_y$ 唯一——映射 $y \mapsto z_y =:
u^*(y)$ 的线性性由唯一性得出。矩阵的对应：
$\langle u(e_i), e_j\rangle$ 按两种方式读出。
\end{proof}

\begin{example}[伴随依赖于内积]\label{ex:b2:quadratic:weightedadjoint}
在 $\R^2$ 上取\emph{加权}内积 $\langle x,
y\rangle_D = x_1y_1 + 2x_2y_2$（矩阵 $D =
\operatorname{diag}(1,2)$），$u$ 在标准基下的矩阵为 $A =
\begin{pmatrix} 0 & 1\\ 0 & 0\end{pmatrix}$。
由 $\langle u(x), y\rangle_D = (Ax)^{\mathsf T}Dy =
x^{\mathsf T}(A^{\mathsf T}D)y$，而 $\langle x, u^*(y)\rangle_D
= x^{\mathsf T}(DA^*)y$，于是\omterm{def:b2:quadratic:adjoint}{伴随}的矩阵为
\[
A^* = D^{-1}A^{\mathsf T}D
= \begin{pmatrix} 1 & 0\\ 0 & \tfrac12\end{pmatrix}
\begin{pmatrix} 0 & 0\\ 1 & 0\end{pmatrix}
\begin{pmatrix} 1 & 0\\ 0 & 2\end{pmatrix}
= \begin{pmatrix} 0 & 0\\ \tfrac12 & 0\end{pmatrix}
\neq A^{\mathsf T} .
\]
验证：取 $x = (1,0)$，$y = (0,1)$：
\[
\langle u(x), y\rangle_D = \langle (0,0), y\rangle_D = 0,
\quad
\langle x, u^*(y)\rangle_D
= \langle(1,0), (0,\tfrac12)\rangle_D = 0 ;
\]
取 $x = (0,1)$，$y = (1,0)$：
\[
\langle u(x), y\rangle_D = \langle(1,0),(1,0)\rangle_D = 1,
\quad
\langle x, u^*(y)\rangle_D =
\langle(0,1),(0,\tfrac12)\rangle_D = 1 .
\] 要点：“$\operatorname{Mat}(u^*)
= \operatorname{Mat}(u)^{\mathsf T}$”只是关于
\emph{标准正交}基的论断；一般而言，度量 $D$
起作用，正如周末问题中的
同时化简那样。
\end{example}

\begin{theorem}[谱定理]\label{thm:b2:quadratic:spectral}
设 $u$ 为欧氏空间 $E$ 上的\omterm{def:b2:quadratic:adjoint}{对称自同态}，则
$E$ 有\emph{由 $u$ 的\omterm{def:b2:reduction:eigen}{特征向量}构成的标准正交基}；所有
\omterm{def:b2:reduction:eigen}{特征值}均为实数，且不同\omterm{def:b2:reduction:eigen}{特征值}的\omterm{def:b2:reduction:eigen}{特征子空间}
相互正交。矩阵形式：每个实对称矩阵 $A$ 可写为
\[
A = P\,D\,P^{\mathsf T},
\qquad P \text{ 正交 } (P^{\mathsf T} P = I),\ D
\text{ 为对角阵 }.\index{谱定理}
\]
\end{theorem}

\begin{proof}
\emph{存在\omterm{def:b2:reduction:eigen}{特征向量}。}函数 $x \mapsto \langle u(x),
x\rangle$ 在单位球面 $S$（$E$ 的）上\omterm{def:b2:metric:continuity}{连续}，该球面
是紧的（有限维，见\cref{thm:b2:nvs:finitedim}），于是
其最大值 $\lambda$ 在某点 $a \in S$ 处取到。断言 $u(a) = \lambda
a$。对任意满足 $y \perp a$、$\norm y = 1$ 的向量与任意 $t \in \R$，
则 $x_t = \frac{a + ty}{\sqrt{1 + t^2}}$ 在 $S$ 上
（$\norm{a + ty}^2 = 1 + t^2$，由勾股定理）；展开
被最大化的函数，
\[
g(t) = \langle u(x_t), x_t\rangle
= \frac{\langle u(a), a\rangle + 2t\langle u(a), y\rangle +
t^2\langle u(y), y\rangle}{1 + t^2}
\]
（$u$ 的对称性合并了两个交叉项：$\langle u(a),
y\rangle = \langle a, u(y)\rangle = \langle u(y), a\rangle$）。
$g$ 是关于 $t$ 的可微函数，且在 $t =
0$ 处取最大值。由商的求导法则，在 $0$ 处给出
\[
g'(0) = \frac{2\langle u(a), y\rangle\cdot 1 - \langle u(a),
a\rangle\cdot 0}{1} = 2\langle u(a), y\rangle = 0 .
\]
于是 $u(a)$ 正交于整个超平面 $a^\perp$：$u(a)
\in (a^{\perp})^{\perp} = \R a$，即 $u(a) = \mu a$；且 $\mu =
\langle u(a), a\rangle = \lambda$。

\emph{归纳法。}正交补空间 $F = a^\perp$ 是
$u$ 不变的：对 $x \perp a$，$\langle u(x), a\rangle = \langle x,
u(a)\rangle = \lambda\langle x, a\rangle = 0$。
限制 $u|_F$ 关于诱导内积是对称的；对
维数归纳，$F$ 有标准正交特征基；并在其前添上 $a$。

\emph{补充。}\omterm{def:b2:reduction:eigen}{特征值}即特征基上的实数 $\langle u(e),
e\rangle$。\omterm{def:b2:reduction:eigen}{特征子空间}的正交性：$u(x) =
\lambda x$，$u(y) = \mu y$ 给出 $\lambda\langle x, y\rangle =
\langle u(x), y\rangle = \langle x, u(y)\rangle = \mu\langle x,
y\rangle$，故 $\langle x, y\rangle = 0$，当 $\lambda \neq \mu$ 时。
矩阵形式：$P$ 的各列即标准正交特征基。
\end{proof}

\begin{example}[完整的谱计算]\label{ex:b2:quadratic:fullrun}
正交对角化 $A = \begin{pmatrix} 1 & 2\\ 2 &
1\end{pmatrix}$。\omterm{def:b2:reduction:charpoly}{特征多项式} $(1 - \lambda)^2 -
4$：\omterm{def:b2:reduction:eigen}{特征值}为 $3$ 与 $-1$。\omterm{def:b2:reduction:eigen}{特征向量}：由 $(A - 3I)v = 0$
得 $v_1 = \frac{1}{\sqrt2}(1,1)$；由 $(A + I)v = 0$ 得 $v_2
= \frac{1}{\sqrt2}(1,-1)$——两者正交，正如
\cref{thm:b2:quadratic:spectral} 所保证的，无需
计算。取 $P = (v_1\ v_2)$（即旋转
$\frac\pi4$）：
\[
P^{\mathsf T}AP = \begin{pmatrix} 3 & 0\\ 0 & -1
\end{pmatrix},
\qquad
x^2 + 4xy + y^2 = 3u^2 - v^2
\quad\text{在旋转后的坐标系中} .
\]
故 \cref{exo:b2:quadratic:1} 中的\omterm{def:b2:quadratic:def}{二次型}属双曲型，
\omterm{thm:b2:quadratic:sylvester}{符号差}为 $(1,1)$，与其高斯配方
$(x + 2y)^2 - 3y^2$ 相符——平方项不同，\omterm{thm:b2:quadratic:sylvester}{符号差}相同，
正如西尔维斯特惯性定律所要求的。小结：高斯法更快给出答案，
但谱方法还表明，在单位
圆上 $q$ 恰好取遍 $\intcc{-1}{3}$，并沿
$v_2$ 与 $v_1$ 取到：额外的计算换来了几何。
\end{example}

\begin{corollary}[主轴；正定性判别]\label{cor:b2:quadratic:principalaxes}
\begin{enumerate}
  \item 欧氏空间上的任一\omterm{def:b2:quadratic:def}{二次型} $q$ 可在某组
        \emph{标准正交}基下对角化：$q(x) = \sum_i \lambda_i
        x_i^2$，其中 $\lambda_i$ 为 $q$ 的对称
        矩阵的\omterm{def:b2:reduction:eigen}{特征值}；\omterm{thm:b2:quadratic:sylvester}{符号差}由正、负
        \omterm{def:b2:reduction:eigen}{特征值}的个数确定。
  \item 对称矩阵半正定（相应地，
        正定）当且仅当其全部\omterm{def:b2:reduction:eigen}{特征值} $\geq 0$（相应地 $>
        0$）；此时瑞利商的
        极值为
        \[
        \min_{\norm x = 1} \langle Ax, x\rangle = \lambda_{\min},
        \qquad
        \max_{\norm x = 1} \langle Ax, x\rangle = \lambda_{\max} .
        \]
\end{enumerate}
\end{corollary}

\begin{proof}
(1) 记 $q(x) = \langle A x, x\rangle$，其中 $A$ 为对称矩阵（
$q$ 在标准正交基下的矩阵）；以谱定理对 $A$
作对角化：对 $x = \sum x_ie_i$，在标准正交
特征基中，
\[
q(x) = \Bigl\langle \sum_i \lambda_ix_ie_i,\ \sum_j
x_je_j\Bigr\rangle = \sum_i \lambda_i x_i^2
\]
（标准正交性使交叉项消失）。$\pm$ 号按
$\lambda_i$ 计数，由西尔维斯特定理给出\omterm{thm:b2:quadratic:sylvester}{符号差}：对每个
坐标乘以 $\sqrt{\abs{\lambda_i}}$ 即得一个高斯
配方，平方项彼此独立。

(2) 在特征基下，$\langle Ax, x\rangle = \sum \lambda_i
x_i^2$，介于 $\lambda_{\min}\norm x^2$ 与
$\lambda_{\max}\norm x^2$ 之间，等号在对应的
\omterm{def:b2:reduction:eigen}{特征向量}处取得；于是全部\omterm{def:b2:reduction:eigen}{特征值}为正，等价于
该\omterm{def:b2:quadratic:def}{二次型}为正定。
\end{proof}

\begin{example}\label{ex:b2:quadratic:spectralexample}
$A = \begin{pmatrix} 2 & 1\\ 1 & 2 \end{pmatrix}$：\omterm{def:b2:reduction:eigen}{特征值} $3$
（\omterm{def:b2:reduction:eigen}{特征向量} $\frac{1}{\sqrt2}(1,1)$）与 $1$
（$\frac{1}{\sqrt2}(1,-1)$）。\omterm{def:b2:quadratic:def}{二次型} $2x^2 + 2xy + 2y^2$
化为 $3X^2 + Y^2$，即在旋转后的标准正交坐标系中：椭圆的
主轴，已算出。高斯配方也能得到对角
形式，但只有谱定理能\emph{在不扭曲
长度}的前提下达到它。
\end{example}

\begin{example}[完整识别一个椭圆]\label{ex:b2:quadratic:ellipserun}
$5x^2 + 4xy + 2y^2 = 6$ 表示什么曲线？其矩阵
$\begin{pmatrix} 5 & 2\\ 2 & 2\end{pmatrix}$ 的\omterm{def:b2:reduction:charpoly}{特征
多项式}为 $\lambda^2 - 7\lambda + 6 = (\lambda - 1)(\lambda -
6)$：\omterm{def:b2:reduction:eigen}{特征值}为 $1$ 与 $6$，均为正——是椭圆。
标准正交\omterm{def:b2:reduction:eigen}{特征向量}：对 $\lambda = 1$，解
$\begin{pmatrix} 4 & 2\\ 2 & 1\end{pmatrix}v = 0$：$v_1 =
\frac{1}{\sqrt5}(1, -2)$；对 $\lambda = 6$：$v_2 =
\frac{1}{\sqrt5}(2, 1)$。在旋转后的坐标 $(X, Y)$
（沿 $(v_2, v_1)$ 方向）中，方程化为
\[
6X^2 + Y^2 = 6,
\qquad\text{即}\qquad
X^2 + \frac{Y^2}{6} = 1 :
\]
半轴 $1$（沿 $v_2$）与 $\sqrt6$（沿 $v_1$）。
小结：粗略形状无需额外计算——$\det
= 6 > 0$ 与迹为正，在任何
\omterm{def:b2:reduction:eigen}{特征向量}算出之前即可宣告为椭圆——但只有谱定理
才给出主轴的方向与长度，即真正的
几何。
\end{example}

\begin{example}[球面上的极值，由谱直接读出]\label{ex:b2:quadratic:sphereextremes}
求 $q(x,y,z) = 2xy + 2yz + 2zx$ 在
单位球面上的极值？其矩阵（即非对角元全为 1 的矩阵，见
\cref{exo:b2:quadratic:2}）的\omterm{def:b2:reduction:eigen}{特征值}为 $2$ 与 $-1$
（二重），由\cref{cor:b2:quadratic:principalaxes}~(2)可知：
\[
\max_{\norm v = 1} q(v) = 2
\ \text{ 在 } v = \tfrac{1}{\sqrt3}(1,1,1),
\qquad
\min_{\norm v = 1} q(v) = -1
\ \text{ 位于圆周 } x + y + z = 0 .
\]
无需微积分，也无需拉格朗日乘数：谱定理
直接解决这一带约束的最优化问题——并给出
极大值点。结尾的体会：可与微分学一章的乘数法对照，
后者同样能求出这些临界点，但计算更繁；
对于球面上的\emph{二次}目标函数，
谱是名副其实的“捷径”（埃尔米特一章的周末问题
正是基于这一观察，建立了整个
柯朗–费希尔理论）。
\end{example}

\begin{remark}[常见误区]
\emph{(i) \omterm{def:b2:quadratic:def}{合同}不等于相似：}\omterm{def:b2:quadratic:def}{二次型}的基变换
作用为 $P^{\mathsf T}BP$，而非 $P^{-1}BP$；\omterm{def:b2:reduction:eigen}{特征值}
并\emph{不是}\omterm{def:b2:quadratic:def}{二次型}的不变量（$I$ 与 $4I$
经 $P = 2I$ 而\omterm{def:b2:quadratic:def}{合同}），只有其符号是
不变的（西尔维斯特）。只有在取定内积之后，
才谈\omterm{def:b2:quadratic:def}{二次型}的\omterm{def:b2:reduction:eigen}{特征值}。\emph{(ii) 元素全为正并不能说明
问题：}$\begin{pmatrix} 1 & 2\\ 2 & 1\end{pmatrix}$ 的元素全为正，
但\omterm{thm:b2:quadratic:sylvester}{符号差}为 $(1,1)$（$\det = -3$）；
反之，正定矩阵的非对角元可能为负，
（将\cref{ex:b2:quadratic:tworoads}平移：
$2I - N$ 同样可行）。可用
\cref{met:b2:quadratic:signature}。\emph{(iii) 平方项
相关：}若写 $q = \ell_1^2 - \ell_2^2$，则
当 $\ell_1, \ell_2$ 成比例时毫无说服力——$x^2 + 2xy + y^2 =
(x+y)^2$ 的秩为 $1$，而非 $2$；务必先检验独立性，
再读取\omterm{thm:b2:quadratic:sylvester}{符号差}。\emph{(iv) 无紧性时的
球面极值：}由
\cref{cor:b2:quadratic:principalaxes}所给的瑞利界均可取到，因为
球面是紧的；在开球或整个空间上，
不定型既无极大值，也无极小值。
\end{remark}

\begin{remark}[应用]
谱定理是本书中应用最广的一条结果：
统计学借助它将协方差矩阵对角化，
（主成分分析），力学由它导出振动的简正模态
（即本章周末问题中的同时化简），
数值分析则以它为基础建立\omterm{pb:b2:quadratic:1}{楚列斯基分解}
与奇异值分解（同一问题），而下一章将其
推广到复埃尔米特空间。本科三年级
分册证明了它在无穷维情形的化身（对于紧的
自伴算子），其中有限维证明里的紧性论证
便成了全部关键所在。
\end{remark}

\begin{remark}[本分册内的展望]
\omterm{def:b2:quadratic:def}{二次型}以三种形态贯穿第 4 册的其余部分：
其一为\emph{海森矩阵}：微分学一章借助
二阶型的\omterm{thm:b2:quadratic:sylvester}{符号差}对临界点进行分类，
因此西尔维斯特不变性使
“鞍点”成为良定义的概念。其二为\emph{能量}：微分方程
一章的振子携带
二次能量 $\frac12x'^2 + \frac12\omega^2x^2$，而
本章周末问题中的同时化简恰是
提取简正模态。其三为\emph{几何}：
本章的圆锥曲线发展为几何各章中的
二次曲面，其中曲面的第二基本形式——
即每个切平面上的\omterm{def:b2:quadratic:def}{二次型}——的\omterm{thm:b2:quadratic:sylvester}{符号差}
决定曲面是像碗那样弯曲，还是
像马鞍那样弯曲。紧接着的埃尔米特一章将整套内容
在 $\C$ 上重演一遍。
\end{remark}

\section{练习}

\begin{exercise}[$\star$]\label{exo:b2:quadratic:1}
用高斯配方法化简，并给出秩与\omterm{thm:b2:quadratic:sylvester}{符号差}：
\[
q_1(x,y) = x^2 + 4xy + y^2,
\qquad
q_2(x,y,z) = x^2 + 2y^2 + 3z^2 + 2xy + 2yz .
\]
\end{exercise}

\begin{exercise}[$\star$]\label{exo:b2:quadratic:2}
正交对角化 $A = \begin{pmatrix} 0 & 1 & 1\\ 1 & 0 &
1\\ 1 & 1 & 0\end{pmatrix}$（\omterm{def:b2:reduction:eigen}{特征值}取自
\cref{ch:b2:reduction}中的计算；现在把基
化为标准正交基），并将\omterm{def:b2:quadratic:def}{二次型} $q(x,y,z) = 2xy + 2yz + 2zx$ 化为
主轴形式。
\end{exercise}

\begin{exercise}[$\star$]\label{exo:b2:quadratic:3}
证明 $u^{**} = u$，$(u \circ v)^* = v^* \circ u^*$，并且
$\ker u^* = (\operatorname{im} u)^{\perp}$。由此推出
$\operatorname{rk} u^* = \operatorname{rk} u$。
\end{exercise}

\begin{exercise}[$\star\star$]\label{exo:b2:quadratic:4}
设 $A$ 为实对称矩阵，且 $A^3 = A$。证明 $A^2$ 为
某个正交投影的矩阵。更一般地，讨论
$A$ 与 $P(A)$ 的谱分解之间的关系（$P$ 为多项式）。
\end{exercise}

\begin{exercise}[$\star\star$]\label{exo:b2:quadratic:5}
证明 $O(n) = \{P : P^{\mathsf T}P = I\}$ 是
$\mathcal{M}_n(\R)$ 的紧子集\emph{（闭：$\{I\}$ 在
\omterm{def:b2:metric:continuity}{连续}映射下的原像；有界：各列为单位向量）}。是否
连通？
\end{exercise}

\begin{exercise}[$\star\star$]\label{exo:b2:quadratic:6}
（平方根）设 $A$ 为对称半正定矩阵，构造
一个对称半正定矩阵 $B$，使 $B^2 = A$，并证明
它是\emph{唯一的}\emph{（存在性：在谱基下对
\omterm{def:b2:reduction:eigen}{特征值}开平方；唯一性：若 $B$ 为候选者，
则其与 $A = B^2$ 可交换，从而保持其各\omterm{def:b2:reduction:eigen}{特征子空间}——归结为
每个子空间上的数量情形）}。
\end{exercise}

\begin{exercise}[$\star\star$]\label{exo:b2:quadratic:7}
对实对称矩阵 $A$，证明 $\vertiii{A}_2 :=
\sup_{\norm x_2 = 1}\norm{Ax}_2 = \max_i \abs{\lambda_i}$
（谱半径），并计算 $\vertiii{A}_2$，其中 $A =
\begin{pmatrix} 1 & 2\\ 2 & 1\end{pmatrix}$.
\end{exercise}

\begin{exercise}[$\star\star\star$]\label{exo:b2:quadratic:8}
（西尔维斯特判别法）设 $A$ 为实对称矩阵，其
顺序主子式 $\Delta_1, \dots, \Delta_n$（左上角
子块的\omterm{def:b2:linalg:det}{行列式}）。证明 $A$ 正定当且仅当
所有 $\Delta_k > 0$。\emph{（对 $\Rightarrow$：正定\omterm{def:b2:quadratic:def}{二次型}的限制
仍为正定，且正定\omterm{def:b2:linalg:det}{矩阵的行列式}——
即其\omterm{def:b2:reduction:eigen}{特征值}之积——为正。
对 $\Leftarrow$：对 $n$ 归纳；左上角 $(n-1)$ 阶子块
为正定，在该子空间上将该\omterm{def:b2:quadratic:def}{二次型}对角化，并
对最后一个变量配方；最后一个对角元的符号
由 $\det A = \Delta_n > 0$ 决定。）}
\end{exercise}

\begin{exercise}[$\star\star\star$]\label{exo:b2:quadratic:9}
（柯朗–费希尔，第二\omterm{def:b2:reduction:eigen}{特征值}）设 $u$ 为\omterm{def:b2:quadratic:adjoint}{对称自同态}，
其\omterm{def:b2:reduction:eigen}{特征值}为 $\lambda_1 \geq \lambda_2 \geq \dots \geq \lambda_n$。
证明
\[
\lambda_2 = \min_{\substack{H \text{超平面}}}\;
\max_{\substack{x \in H,\ \norm x = 1}} \langle u(x), x\rangle .
\]
\emph{（对 $\leq$：任一超平面与由
前两个\omterm{def:b2:reduction:eigen}{特征向量}张成的 $2$ 维平面相交。对 $\geq$：取 $H =
(e_1)^{\perp}$.)}
\end{exercise}

\begin{exercise}[$\star\star$]\label{exo:b2:quadratic:10}
确定 $q(x_1, \dots, x_n) = \sum_{i <
j} x_ix_j$ 在 $\R^n$ 上的秩与\omterm{thm:b2:quadratic:sylvester}{符号差}（$n \geq 2$），用两种方法求：先由代数恒等式
$2q = \bigl(\sum x_i\bigr)^2 - \sum x_i^2$ 并结合
$q$ 在超平面 $\sum x_i = 0$ 上的限制；
另一种是计算其矩阵 $\frac12(J - I)$ 的\omterm{def:b2:reduction:eigen}{特征值}，
其中 $J$ 为全 1 矩阵。
\end{exercise}

\begin{exercise}[$\star\star$]\label{exo:b2:quadratic:11}
设 $A$、$B$ 为实对称矩阵，且 $B$ 为半正定矩阵。
证明
\[
\lambda_{\min}(A)\operatorname{tr} B
\;\leq\;
\operatorname{tr}(AB)
\;\leq\;
\lambda_{\max}(A)\operatorname{tr} B .
\]
\emph{（记 $B = C^{\mathsf T}C$，且 $\operatorname{tr}(AB) =
\sum_i \langle A c_i, c_i\rangle$，其中 $c_i$ 取遍
$C^{\mathsf T}$ 的各列）}。特别地，
当两者均为半正定矩阵时，$\operatorname{tr}(AB) \geq 0$。
\end{exercise}

\begin{exercise}[$\star\star\star$]\label{exo:b2:quadratic:12}
在 $E = \mathcal{M}_n(\R)$ 上考虑 $q(M) =
\operatorname{tr}(M^2)$.
\begin{enumerate}
  \item 证明 $q$ 是\omterm{def:b2:quadratic:def}{二次型}，其极化双线性型为
        $\varphi(M, N) = \operatorname{tr}(MN)$.
  \item 证明对称矩阵与反对称矩阵构成
        关于 $\varphi$ 正交的子空间，$q$ 在其上分别
        为正定与负定
        \emph{（对各情形分别逐元计算 $\operatorname{tr}(M^2)$，
        即可验证）}。
  \item 结论：$q$ 的\omterm{thm:b2:quadratic:sylvester}{符号差}为 $\bigl(\frac{n(n+1)}{2},
        \frac{n(n-1)}{2}\bigr)$，秩为 $n^2$。
\end{enumerate}
\end{exercise}

\section{问题：楚列斯基、阿达马与
极分解}

\begin{problem}\label{pb:b2:quadratic:1}
谱定理是一台显微镜；本题则将它用作
一座工厂。由格拉姆矩阵出发，我们构造\emph{楚列斯基
分解}\index{楚列斯基分解}（并
将高斯主元识别为子式之比），然后证明\emph{阿达马
不等式}\index{阿达马不等式}（关于\omterm{def:b2:linalg:det}{行列式}），并构造
\emph{极分解}\index{极分解} $A =
QS$ 及\emph{奇异值分解}，将平面
二次曲线分类，最后以两个\omterm{def:b2:quadratic:def}{二次型}的同时化简——
即振动简正模式背后的定理。全文中，$E =
\R^n$ 配以其标准内积。

\medskip
\noindent\textbf{第一部分——格拉姆矩阵与\omterm{pb:b2:quadratic:1}{楚列斯基分解}。} 对于
向量 $v_1, \dots, v_n \in E$，其\emph{格拉姆矩阵}是
$G = \bigl(\langle v_i, v_j\rangle\bigr)_{i,j}$.
\begin{enumerate}
  \item 证明 $G$ 为对称半正定矩阵，
        并且当且仅当 $(v_1, \dots, v_n)$
        线性无关时为正定 \emph{（计算 $X^{\mathsf T}GX$）}。
  \item 反之，证明每个对称
        半正定矩阵 $A$ 都是格拉姆矩阵：$A = C^{\mathsf T}C$，
        对某个 $C$（利用
        \cref{exo:b2:quadratic:6}中的平方根），且 $C$ 可逆当且仅当 $A$
        为正定。
  \item 由此推出，半正定矩阵 $A$ 满足
        $\abs{a_{ij}} \leq \sqrt{a_{ii}\,a_{jj}}$ 对所有 $i,
        j$ \emph{（限制到两个坐标上）}——这正是
        以矩阵语言重读的柯西–施瓦茨不等式。
  \item （\omterm{pb:b2:quadratic:1}{楚列斯基分解}）设 $A$ 为正定矩阵，证明
        存在\emph{唯一的}上三角矩阵 $T$，其
        对角元均为正，使得
        \[
        A = T^{\mathsf T}\,T
        \]
        \emph{（存在性：对使
        $A$ 成为格拉姆矩阵的向量施以格拉姆–施密特正交化；唯一性：若
        $T_1^{\mathsf T}T_1 = T_2^{\mathsf T}T_2$，证明
        $T_1T_2^{-1}$ 为正交上三角矩阵，其
        对角元为正，从而为 $I$）}。
  \item 证明顺序主子式满足
        $\Delta_k = (t_{11}\cdots t_{kk})^2$，并由此推出
        正定\omterm{def:b2:quadratic:def}{二次型}经高斯消元所产生的主元，
        按自然的变量
        次序取定，依次为
        \[
        d_k = \frac{\Delta_k}{\Delta_{k-1}}
        \qquad (\Delta_0 = 1) :
        \]
        西尔维斯特判别法中的子式
        （\cref{exo:b2:quadratic:8}）与高斯主元
        是同一组数据。在
        \cref{ex:b2:quadratic:tworoads}上验证。
\end{enumerate}

\medskip
\noindent\textbf{第二部分——\omterm{pb:b2:quadratic:1}{阿达马不等式}。}
\begin{enumerate}[resume]
  \item 设 $A$ 为正定矩阵。证明
        \[
        \det A \leq a_{11}\,a_{22}\cdots a_{nn}
        \]
        \emph{（归一化：令 $B = DAD$，其中 $D =
        \operatorname{diag}(a_{ii}^{-1/2})$ 使对角元全为 1；
        用算术–几何平均不等式估计 $\det B = \prod \mu_i$，
        与 $\operatorname{tr} B = n$ 比较）}。
  \item 证明等号成立当且仅当 $A$ 为对角矩阵。
  \item 由此推出\emph{\omterm{pb:b2:quadratic:1}{阿达马不等式}}：对每个实
        方阵 $M$，其列为 $c_1, \dots, c_n$，
        \[
        \abs{\det M} \leq \prod_{i=1}^{n}\norm{c_i}_2 ,
        \]
        当 $M$ 可逆时，等号成立当且仅当
        各列两两正交\emph{（将第 6--7 题用于
        $M^{\mathsf T}M$）}。
  \item 几何与组合方面的推论：将第 8 题理解为
        “平行多面体的体积
        不超过其各棱长之积”；并证明
        元素均在 $\intcc{-1}{1}$ 内的矩阵满足
        $\abs{\det M} \leq n^{n/2}$。（达到此界的矩阵——阿达马矩阵——
        存在于 $n = 1, 2$
        及许多 $4$ 的倍数；是否对\emph{所有}
        $4$ 的倍数都存在，是一个著名的公开问题）。
\end{enumerate}

\medskip
\noindent\textbf{第三部分——\omterm{pb:b2:quadratic:1}{极分解}与
奇异值。}
\begin{enumerate}[resume]
  \item 设 $A$ 可逆。证明 $A^{\mathsf T}A$
        为正定矩阵，并且
        \[
        S = \sqrt{A^{\mathsf T}A}
        \quad\text{（
        \cref{exo:b2:quadratic:6}的平方根）},
        \qquad Q = AS^{-1}
        \]
        于是给出分解 $A = QS$，其中 $Q$ 为正交矩阵，
        $S$ 为正定矩阵。
  \item 证明可逆矩阵 $A$ 的上述分解
        是唯一的。
  \item 将存在性推广到任意 $A$：取
        $\varepsilon_k \to 0$，使 $A + \varepsilon_k I$
        可逆，写出 $A + \varepsilon_kI = Q_kS_k$，
        并利用 $O(n)$ 的紧性
        （见\cref{exo:b2:quadratic:5}），取子列使 $Q_k \to Q$；
        证明 $S_k = Q_k^{\mathsf T}(A + \varepsilon_kI)$
        收敛到某个半正定矩阵 $S$，使得 $A =
        QS$，且 $S = \sqrt{A^{\mathsf T}A}$。
        对奇异的 $A$，唯一性在何处失效？
  \item （奇异值分解）由此可知，每个实
        方阵 $A$ 可写成
        \[
        A = U\,\Sigma\,V^{\mathsf T},
        \qquad U, V \in O(n),\quad
        \Sigma = \operatorname{diag}(\sigma_1, \dots,
        \sigma_n),\ \sigma_i \geq 0 ,
        \]
        其中 $\sigma_i$（即\emph{奇异值}）
        为 $\sqrt{A^{\mathsf T}A}$ 的\omterm{def:b2:reduction:eigen}{特征值}。
  \item 三个推论：$\vertiii{A}_2 = \sigma_{\max}$ 对
        \emph{每个}实矩阵 $A$ 成立（推广
        \cref{exo:b2:quadratic:7}）；$\abs{\det A} =
        \sigma_1\cdots\sigma_n$；以及单位球面
        在可逆矩阵 $A$ 作用下的像为椭球面，
        其半轴长分别为 $\sigma_1, \dots, \sigma_n$，沿
        $U$ 的各列方向。
\end{enumerate}

\medskip
\noindent\textbf{第四部分——借助谱定理讨论二次曲线。}
平面上的二次曲线是函数 $f(x) = q(x) + \langle b,
x\rangle + c$ 的零点集，其中 $q \neq 0$ 是矩阵为 $A$ 的\omterm{def:b2:quadratic:def}{二次型}，
$b \in \R^2$，$c \in \R$。
\begin{enumerate}[resume]
  \item 对 $f$ 先作旋转（主轴化，
        见\cref{cor:b2:quadratic:principalaxes}），再作
        平移，并对可能的非空
        非退化图形按 $q$ 的\omterm{thm:b2:quadratic:sylvester}{符号差}分类：椭圆
        （$\det A > 0$）、双曲线（$\det A < 0$）、抛物线（$\det A = 0$，
        秩为 $1$，且一次项未被吸收
        ）。
  \item 对下列方程完整地作化简：
        \[
        x^2 + 4xy + y^2 + 2x - 2y = 4 :
        \]
        旋转后的坐标、化简后的方程、曲线的类型与
        中心。
  \item （中心二次曲线）设 $\det A \neq 0$。证明
        \emph{中心}为 $x_0 = -\frac12 A^{-1}b$，并
        以 $\begin{pmatrix} I & x_0\\ 0 &
        1\end{pmatrix}$ 对 $3\times3$ 矩阵 $\widetilde Q
        = \begin{pmatrix} A & b/2 \\ b^{\mathsf T}/2 &
        c\end{pmatrix}$ 作\omterm{def:b2:quadratic:def}{合同}变换，得到
        \[
        \det\widetilde Q = f(x_0)\,\det A :
        \]
        中心二次曲线退化（为一点或两条直线）
        当且仅当 $\det\widetilde Q = 0$ 时。
  \item 在第 16 题的例子上验证第 17 题：
        计算 $x_0$、$f(x_0)$ 与 $\det\widetilde Q$，并
        再次断定该二次曲线为非退化
        双曲线。
  \item （二次曲面束）对 $\lambda \in \R$，分类
        曲面
        \[
        x^2 + y^2 + z^2 + 2\lambda(xy + yz + zx) = 1
        \]
        依其矩阵的\omterm{def:b2:reduction:eigen}{特征值}，\emph{（全 1 结构：
        \omterm{def:b2:reduction:eigen}{特征值} $1 + 2\lambda$ 与 $1 -
        \lambda$ 为二重）}：球面或椭球面、柱面、一对
        平面、单叶与双叶双曲面，依
        $\lambda$ 而定。
\end{enumerate}

\medskip
\noindent\textbf{第五部分——同时处理两个\omterm{def:b2:quadratic:def}{二次型}：
同时化简。}
\begin{enumerate}[resume]
  \item 设 $q$ 为正定的，$q'$ 为
        $E$ 上任意的\omterm{def:b2:quadratic:def}{二次型}。证明
        存在 $E$ 的一组基，它对 $q$ 标准正交，
        又对 $q'$ 正交：在该基下，$q = \sum x_i^2$ 且 $q' = \sum \mu_i
        x_i^2$ \emph{（取 $q$ 为内积，并将
        谱定理应用于表示
        $q'$ 的自同态）}。
  \item 矩阵形式：若 $A$ 正定、$B$ 对称，则
        存在可逆矩阵 $P$，使得
        $P^{\mathsf T}AP = I$ 且 $P^{\mathsf T}BP =
        \operatorname{diag}(\mu_1, \dots, \mu_n)$，其中
        $\mu_i$ 为 $\det(B - \mu A) = 0$ 的根。
  \item 对下列矩阵完整地计算
        \[
        A = \begin{pmatrix} 2 & 1\\ 1 & 1 \end{pmatrix},
        \qquad
        B = \begin{pmatrix} 0 & 1\\ 1 & 0 \end{pmatrix} :
        \]
        求广义\omterm{def:b2:reduction:eigen}{特征值} $\mu_\pm$，以及
        使两个\omterm{def:b2:quadratic:def}{二次型}同时对角化的向量。
  \item 说明正定性不能省去：对于
        \[
        A = \begin{pmatrix} 1 & 0\\ 0 & -1 \end{pmatrix},
        \qquad
        B = \begin{pmatrix} 0 & 1\\ 1 & 0 \end{pmatrix},
        \]
        不存在能同时对角化两个\omterm{def:b2:quadratic:def}{二次型}的基\emph{（若 $P$
        同时对角化两者，则 $\det(B - \mu A)$ 将
        分解为实根之积；请计算它）}。
  \item 说明第 21 题中的 $\mu_i$ 即为
        $A^{-1}B$ 的\omterm{def:b2:reduction:eigen}{特征值}，并说明 $A^{-1}B$ 虽
        一般不是对称阵，但总\omterm{def:b2:reduction:diag}{可对角化}，
        且\omterm{def:b2:reduction:eigen}{特征值}为实数\emph{（用 $\sqrt A$ 作相似变换）}。
  \item 综述。各用一句话回答：(i) 各部分共同依赖的那一条定理；
        (ii) 第一至三部分中哪些结论对
        正\emph{半}定矩阵仍然成立；
        哪些需要正定性；(iii) 物理系统
        的微小振动为第 20--22 题所对角化的
        系统（动能与势能分别作为两个\omterm{def:b2:quadratic:def}{二次型}），
        以及 $\mu_i$ 在其中的含义。
\end{enumerate}
\end{problem}
